%% SCRAP: experiments/campaigns/l8_attractor_map/results_20251209_145907/binary_2cycle_analysis_report
%% SOURCE: docs/working/experiments/campaigns/l8_attractor_map/results_20251209_145907/binary_2cycle_analysis_report.md
%% STATUS: HISTORICAL
%% FITS: experiments/ch-l8-map
%% EDITORIAL: lifted — prose rewritten to press voice

\section{Binary 2-Cycle Analysis: Phase-Space Clustering}

\textbf{Date:} 11 December 2025. \textbf{Data source:}
\texttt{heartbeat\_all.csv} from the L8 Attractor Map experiment.

\subsection{Method}

The analysis applied k-means clustering ($k = 2$) to the (HR, $\Delta$HR)
coordinate space, where HR is the tick interval in nanoseconds and
$\Delta$HR is the first difference of successive intervals. A binary
clarity metric quantifies the separation of the two clusters:

\[
  \text{BinaryClarity} = \frac{\text{separation}^2}{\text{avg\_spread}}
\]

where separation is the Euclidean distance between cluster centers and
avg\_spread is the average effective radius of the two covariance ellipses.
Outliers beyond $5\sigma$ were removed before clustering.

\subsection{Results}

\begin{center}
\begin{tabular}{llrrl}
\toprule
Rank & Workload   & Binary Clarity & Separation (ns) & Character \\
\midrule
1 & TRANSITION & $1.27 \times 10^6$ & 258{,}905 & Dominant two-regime structure \\
2 & VOLATILE   & $7.66 \times 10^4$ & 41{,}487  & Moderate separation \\
3 & TEMPORAL   & $7.02 \times 10^4$ & 44{,}851  & Moderate separation \\
4 & OMNI       & $6.17 \times 10^4$ & 42{,}347  & Moderate separation \\
5 & STABLE     & $6.09 \times 10^4$ & 35{,}932  & Weakest separation \\
6 & DIVERSE    & $6.05 \times 10^4$ & 42{,}846  & Moderate separation \\
\bottomrule
\end{tabular}
\end{center}

\subsection{Interpretation}

TRANSITION exhibited binary clarity approximately $16\times$ higher than
the other five workloads. The two cluster centers were separated by
258{,}905\,ns — a fast-heartbeat regime versus a slow-heartbeat regime —
compared to 35{,}000--45{,}000\,ns for the remaining workloads.

The result is counterintuitive: STABLE shows the \emph{weakest} binary
structure despite its name. STABLE operates in a homogeneous regime with
tight cluster overlap, while TRANSITION's explicit state switching creates
the most pronounced phase-space separation. The 2-cycle structure is
present across all workloads; it is most clearly exposed when the workload
forces repeated transitions between distinct operational states.

Loop \#7 (Adaptive Heartrate) is the likely mechanism: it creates attractor
points in phase space by adjusting tick frequency in response to workload
characteristics, producing the binary clustering even in nominally
non-transitioning workloads.

\subsection{Open Questions}

\begin{enumerate}
  \item TRANSITION: are exactly two heartbeat regimes present, or is $k=2$
        an oversimplification? Silhouette scores for $k = 3, 4$ are
        warranted.
  \item Autocorrelation analysis of HR$(t)$ for individual runs would
        detect periodic switching behavior directly.
  \item The Levene's test used internally by Loop \#5 (Window Width
        Inference) should be compared against cluster variance homogeneity
        to test for cross-loop interaction.
\end{enumerate}

\subsection{Generated Artifacts}

\begin{itemize}
  \item \texttt{binary\_2cycle\_analysis.svg} — six-panel visualization
        ($1800 \times 1200$\,px), one panel per workload
  \item \texttt{analyze\_2cycle.py} — pure Python analysis script (no
        external dependencies)
\end{itemize}
